\scanpage{132}
\mspage{-III- \handwritten{4} -}
\begin{flushright}n° \handwritten{383}\end{flushright}
$s_i$, donc pour $G$ (prop.1), et est par conséquent orthogonal à $W$ ; ceci montre que les $e_i$ engendrent $W$. C.Q.F.D.

Ceci étant, soient $I=\handwritten{\{}1,2,\ldots,m\handwritten{\}}$, $M=(m_{ij})$ la matrice de Coxeter sur $I$ où $m_{ij}$ est l'ordre de $s_i s_j$ dans $G$, et $G(M)$ le groupe de Coxeter défini par les générateurs $r_i$ $(i\in I)$ et les relations $(r_i r_j)^{m_{ij}}=1$ (\S\ 2, n°1). D'après la prop.3 on peut prendre $W$, muni de la base $(e_i)$, pour espace de la représentation directe $\handwritten{\rho}$ de $G(M)$ ; d'après la prop.2 et le fait que les $e_i$ sont unitaires, la forme bilinéaire $B_M$ du \S\ 2, n°2 est la restriction à $W$ de la forme donnée $(x\handwritten{|}y)$ ; ainsi $\handwritten{\rho}(r_i)$ est la restriction $s'_i$ à $W$ de la réflexion $s_i$. Identifions $W$ à son dual au moyen de la forme $(x\handwritten{|}y)$ ; comme, pour passer d'une symétrie à sa contragrédiente, on échange vecteur et forme linéaire (\S\ 2, n°3), et comme $s'_i$ est définie par un vecteur unitaire et l'hyperplan orthogonal, $s'_i$ est sa propre contragrédiente. Si $(e'_i)$ est la base de $W$ duale de $(e_i)$ (i.e. $(e_i\handwritten{|}e'_j)=\handwritten{\delta}_{ij}$ pour tout $i,j$), l'ensemble des $\displaystyle x=\handwritten{\sum}_{i\in I}a_i e'_i$ $(a_i\in R)$ de $W$ tels que $(x\handwritten{|}e_i)>0$ pour tout $i$ est celui des $x$ tels que $a_i>0$ pour tout $i$ ; donc, d'après (2), $C\handwritten{\cap}W$ est la chambre notée $C$ au \S\ 2, n°3. En tenant compte du fait que $G$ opère trivialement sur le sous-espace $W'$ orthogonal à $W$, le th.1 du \S\ 2, n°3 donne le résultat suivant :

THEOREME 1.- \uline{Soient} $\underline{H}$ \uline{un ensemble d'hyperplans de} $E$ \uline{vérifiant} (1) \uline{et} (2), $C$ \uline{une chambre définie par} $\underline{H}$, $H_1,\ldots,H_m$ \uline{ses murs}, \uline{et} $s_i$ \uline{la réflexion par rapport à} $H_i$. \uline{Notons} $m_{ij}$ \uline{l'ordre de} $s_i s_j$ \uline{dans} $G$. \uline{Alors}

a) \uline{les hyperplans} $H_i$ \uline{sont linéairement indépendants} ;

b) \uline{le groupe} $G$ \uline{est défini par les générateurs} $s_i$ \uline{et les relations} $s_i^2=(s_i s_j)^{m_{ij}}=1$;

c) \uline{le groupe} $G$ \uline{opère de façon simplement transitive sur l'ensemble des chambres définies par} $\underline{H}$.

COROLLAIRE 1.- \uline{Si} $G$ \uline{est essentiel}, \uline{deux vecteurs quelconques} $x,y$ \uline{d'une même chambre} $C$ \uline{font un angle aigu} (i.e. $(x\handwritten{|}y)\handwritten{\geq}0$).
% archive sheet number: 131
